% =====================================================================
%  corps/01-introduction.tex
% =====================================================================
\chapter{Introduction}
\label{chap:introduction}

Ce chapitre introduit la problématique de recherche, le contexte
scientifique et la structure de la thèse.

\section{Contexte et motivation}

La traductologie contemporaine s'appuie de plus en plus sur des
méthodes outillées et des corpus multilingues~\autocite{toury1995}.
Comme le rappelle \textcite{eco2003}, traduire revient souvent à
« dire presque la même chose ». Les travaux menés au sein du
Département TIM illustrent cette évolution, qu'il s'agisse de
traduction automatique de dialectes~\autocite{alalmaoui2025arabizi,mutal2026aladdin},
de traduction de la parole en milieu médical~\autocite{bouillon2021babeldr},
de terminologie et de variation~\autocite{humbertdroz2022higgs}, de
localisation~\autocite{morado2020xliff} ou encore de détection
automatique du changement sémantique~\autocite{tahmasebi2021semanticchange,schlechtweg2020semeval}.

\section{Question de recherche}

% Formulez ici votre question de recherche et vos hypothèses.
Nous posons la question suivante : \dots

\section{Structure de la thèse}

La thèse est organisée comme suit. Le
chapitre~\ref{chap:exemples} illustre les capacités multilingues du
présent modèle. La conclusion (chapitre~\ref{chap:conclusion})
synthétise les apports.
